\documentclass[12pt]{article}
\usepackage[utf8]{inputenc}
\usepackage[T1]{fontenc}
\usepackage[margin=1in]{geometry}
\usepackage{setspace}
\usepackage{xcolor}
\usepackage[colorlinks=true, urlcolor=blue!60!black, linkcolor=blue!60!black]{hyperref}
\setlength{\parindent}{2em}
\setlength{\parskip}{1.25em}
\renewcommand{\baselinestretch}{1.0}
\begin{document}
\begin{center}
{\LARGE \textit{Shelley v. Kraemer}, 334 U.S. 1 (1948)\par}
\vspace{0.8em}
{\large Case Brief --- Four-Page Edition\par}
\end{center}
\vspace{1em}

\section{Memory Jogger}

Private racially restrictive covenants do not themselves violate the Fourteenth Amendment, but a state court's enforcement of them is state action that denies equal protection; a unanimous six-Justice Court reversed Missouri and Michigan decrees ousting Black families from homes they had bought from willing sellers (\href{https://supreme.justia.com/cases/federal/us/334/1/}{Justia}).

\section{Detailed Case Facts}

\textbf{Missouri.} In 1911, thirty of thirty-nine owners on a stretch of Labadie Avenue in St. Louis recorded a fifty-year covenant providing that no part of their property could be ``occupied by any person not of the Caucasian race.'' Five parcels in the district were already owned by Black residents. In 1945 J.D. and Ethel Shelley, a Black couple, bought one of the restricted parcels through an intermediary, without actual knowledge of the covenant.

\textbf{Michigan.} A 1934 covenant on a Detroit property provided: ``This property shall not be used or occupied by any person or persons except those of the Caucasian race.'' In 1944 Orsel and Minnie McGhee, a Black couple, acquired title and moved in.

\section{Procedural History}

In Missouri, neighbors Louis and Fern Kraemer sued to restrain the Shelleys from taking possession and to divest their title. The trial court denied relief because the covenant had never taken effect, but the Missouri Supreme Court reversed, ordered enforcement, and found no federal violation. In Michigan, neighbors sued and the trial court ordered the McGhees to vacate within ninety days; the Michigan Supreme Court affirmed. The U.S. Supreme Court consolidated the cases, heard argument in January 1948, and decided them on May 3, 1948.

\section{Judicial Votes}

Reversed, \textbf{6--0}.
\begin{itemize}
\item \textbf{Majority:} Chief Justice Vinson, joined by Justices Black, Frankfurter, Douglas, Murphy, and Burton.
\item \textbf{Concurrences and dissents:} None.
\item \textbf{Not participating:} Justices Reed, Jackson, and Rutledge, whose recusals scholarship attributes to their ownership of covenant-restricted homes.
\end{itemize}

\section{Holding}

Private agreements excluding persons of a designated race from real estate, standing alone and effectuated only through voluntary adherence, do not violate the Fourteenth Amendment. But the actions of state courts and judicial officers in their official capacities are state action, even when they enforce private agreements under common-law rules. By enforcing the covenants to divest title or enjoin occupancy by willing, financially able buyers from willing sellers, Missouri and Michigan denied equal protection. The Court did not reach the due process and privileges or immunities claims.

\section{Analysis of Opinions}

\textbf{Precedent and property rights.} \textit{Corrigan v. Buckley} arose in the District of Columbia and addressed only the covenant's validity; \textit{Hansberry v. Lee} turned on procedural due process. The covenants targeted a class defined ``wholly in terms of race or color,'' and equal property rights were a core aim of the Amendment's framers, reflected in the Civil Rights Act of 1866. Had the States imposed these bans by statute, \textit{Buchanan v. Warley} would condemn them. But under the \textit{Civil Rights Cases}, the Amendment reaches only state action.

\textbf{Courts as state actors.} Since \textit{Ex parte Virginia} (``A State acts by its legislative, its executive, or its judicial authorities. It can act in no other way.''), judicial officers have been treated as the State. That holds not only for unfair trials but for enforcement of substantive common-law rules through impeccable proceedings, as \textit{AFL v. Swing} and \textit{Cantwell v. Connecticut} show.

\textbf{Active intervention.} State action existed ``in the full and complete sense of the phrase'': but for the courts' intervention, the families would have occupied their homes freely. The Court rejected the argument that courts would equally enforce covenants against white buyers, because Fourteenth Amendment rights are ``personal rights,'' and held that covenantors have no right to demand state action denying others equal protection.

\textbf{Companion case and aftermath.} In \textit{Hurd v. Hodge}, the Court barred enforcement in the District of Columbia, Frankfurter concurring on equitable grounds. \textit{Barrows v. Jackson} (1953) extended \textit{Shelley} to damages actions, while \textit{Evans v. Abney} (1970) allowed neutral trust-reversion law to close a segregated park.

\section{Examples: Future Applications}

\begin{enumerate}
\item \textbf{Same-Side (Barred): Religious Covenant.} An association sues to divest title after a willing owner sells to a Muslim family; third parties are enlisting judicial coercion to override a voluntary sale.
\item \textbf{Same-Side (Barred): Damages Against the Seller.} Neighbors sue an owner who sold to a buyer of an excluded national origin; under \textit{Barrows}, the threat of liability would pressure owners to refuse willing buyers.
\item \textbf{Opposite-Side (No Violation): Trespass From One's Own Home.} A homeowner who excludes a guest for religious reasons and calls police is enforcing her own choice, not overriding a willing owner.
\item \textbf{Opposite-Side (May Proceed): Non-Disparagement Damages.} Under \textit{Cohen v. Cowles Media}, enforcing a voluntarily assumed obligation under generally applicable law need not violate the First Amendment.
\item \textbf{Fence-Sitter: Faith-Based Bequest Condition.} A will conditions a bequest on the grandson's marrying within the faith. The decree would give legal force to a private religious exclusion, but no willing parties are overridden, as in \textit{Evans v. Abney}.
\end{enumerate}

\section{Critique}

\textbf{Collapse of the public/private line.} Because nearly all property, contract, and probate rights depend on judicial backing, a literal reading of \textit{Shelley} would make every private choice state action; lower courts therefore confine it to racial discrimination, refusing, for example, to treat enforcement of speech-restricting settlements as state action.

\textbf{Narrower reconstructions.} After Herbert Wechsler charged that \textit{Shelley} lacked a neutral principle, scholars proposed willing-party coercion, de facto private zoning, or the Thirteenth Amendment and \S~1982 route later taken in \textit{Jones v. Alfred H. Mayer Co.} Critics also note that in \textit{Swing} the State authored the restriction, whereas here private covenantors did.

\section{Key Quotations}

\begin{quote}
``That Amendment erects no shield against merely private conduct, however discriminatory or wrongful.''
\end{quote}
\noindent --- Chief Justice Vinson, for the Court, 334 U.S. at 13.

\begin{quote}
``It is clear that, but for the active intervention of the state courts, supported by the full panoply of state power, petitioners would have been free to occupy the properties in question without restraint. \ldots\ Rather, these are cases in which the States have made available to such individuals the full coercive power of government to deny to petitioners, on the grounds of race or color, the enjoyment of property rights in premises which petitioners are willing and financially able to acquire and which the grantors are willing to sell.''
\end{quote}
\noindent --- Chief Justice Vinson, for the Court, 334 U.S. at 19.

\begin{quote}
``The rights created by the first section of the Fourteenth Amendment are, by its terms, guaranteed to the individual. The rights established are personal rights. It is, therefore, no answer to these petitioners to say that the courts may also be induced to deny white persons rights of ownership and occupancy on grounds of race or color. Equal protection of the laws is not achieved through indiscriminate imposition of inequalities.''
\end{quote}
\noindent --- Chief Justice Vinson, for the Court, 334 U.S. at 22.

\end{document}
